\metatitle{Ultra-log-concavity for rigid width-three posets}
\metadescription{A three-color matroid encoding and Lean proof of
ultra-log-concavity for P-Eulerian polynomials of rigid width-three posets.}

\section[widthThreeLogConcavity]{Ultra-log-concavity for rigid width-three posets}

The $P$-Eulerian polynomial records descents in the linear extensions of a
poset. Brenti conjectured that its coefficients are log-concave with no
internal zeros \cite{Brenti1989}. Real-rootedness is false even for naturally
labeled width-two posets \cite{Stembridge2007NeggersStanleyCounter}, while
ultra-log-concavity holds in width two by the rook-matroid construction of
Alexandersson and Jal \cite{AlexanderssonJal2024x}.

Here we record a formalized extension to a class of width-three posets. The
proof encodes linear extensions by three-color words whose two hidden color
classes are bases of transversal matroids. Rigidity makes this encoding exact.

\subsection[posetEulerianPolynomials]{$P$-Eulerian polynomials}

Let $P$ be a naturally labeled poset on $[n]=\{1,\dotsc,n\}$. A linear
extension is a permutation $\sigma=\sigma_1\dotsm\sigma_n$ compatible with
the order of $P$. Write $\mathcal L(P)$ for the set of linear extensions. A
position $i<n$ is a descent if $\sigma_i>\sigma_{i+1}$.

\begin{definition}[$P$-Eulerian polynomial]
The \defin{$P$-Eulerian polynomial} is
\[
 W_P(t)\coloneqq
 \sum_{\sigma\in\mathcal L(P)}t^{\operatorname{des}(\sigma)}.
\]
\end{definition}

\begin{definition}[Ultra-log-concavity]
A nonnegative sequence $(a_0,\dotsc,a_d)$ is
\defin{ultra-log-concave of order $d$} if
\[
 \left(\frac{a_j}{\binom d j}\right)^2
 \geq
 \frac{a_{j-1}}{\binom d{j-1}}
 \frac{a_{j+1}}{\binom d{j+1}}
 \qquad (1\leq j<d).
\]
It has \defin{no internal zeros} if $a_i,a_k>0$ and $i<j<k$ imply
$a_j>0$.
\end{definition}

\subsection[rigidChoices]{Rigid choices}

Let $Z=z_1<\dotsm<z_h$ be a longest chain of $P$, and let
$X=x_1<\dotsm<x_r$ and $Y=y_1<\dotsm<y_s$ be chains that partition
$P\setminus Z$. Thus $r+s=n-h$. For $q\notin Z$, put
\[
 p_q\coloneqq\#\{z\in Z:z<q\},
 \qquad
 b_q\coloneqq h-\#\{z\in Z:z>q\}.
\]
The constrained order is generated by the order of $P$ together with
\[
 z<q\quad\text{if }z\in Z\text{ and }z\parallel q,
 \qquad
 y<x\quad\text{if }x\in X,\ y\in Y\text{ and }x\parallel y.
\]

For $x\in X$, define
\[
 \operatorname{lo}(x)\coloneqq\#\{y\in Y:y<x\},
 \qquad
 \operatorname{hi}(x)\coloneqq s-\#\{y\in Y:y>x\}.
\]
If $x=x_i$, let $K_x$ be the set of $k\in\{0,\dotsc,s\}$ such that
\[
 \operatorname{lo}(x)\leq k\leq\operatorname{hi}(x)
\]
and either $k=\operatorname{hi}(x)$, or $i<r$ and
$k\geq\operatorname{lo}(x_{i+1})$. We use the conventions
$p_{y_0}=0$ and $b_{y_{s+1}}=h$.

\begin{definition}[Rigid choice]
The triple $(Z,X,Y)$ is a \defin{rigid choice} if the constrained order is
acyclic and the following conditions hold.
\begin{itemize}
\item \textbf{(D)} If $x\in X$, $y\in Y$, and $x\parallel y$, then
$p_x\leq p_y$.
\item \textbf{(Y)} For each $1\leq k\leq s$, the number
$\max(p_x,p_{y_{k-1}})$ is independent of the choice of
$x\in X$ with $x\parallel y_k$.
\item \textbf{(X)} For each $x\in X$, the pair
\[
 \bigl(\max(p_x,p_{y_k}),\min(b_x,b_{y_{k+1}})\bigr)
\]
is independent of $k\in K_x$.
\end{itemize}
\end{definition}

Since $P$ is covered by the three chains $Z,X,Y$, it has width at most three.

\subsection[fullThreeColorFamilies]{Full three-color matroid families}

For a word $u:E\to\{0,1,2\}$, write
\[
 B(u,c)\coloneqq\{e\in E:u(e)=c\}.
\]
Let $M_1$ and $M_2$ be matroids on $E$. Consider the family
\[
 \mathcal F(M_1,M_2)\coloneqq
 \{u:E\longrightarrow\{0,1,2\}:
 B(u,1)\in\mathcal B(M_1),\ B(u,2)\in\mathcal B(M_2)\}.
\]
The remaining positions have color $0$. Its target quota is
\[
 R=|E|-\operatorname{rk}(M_1)-\operatorname{rk}(M_2),
\]
and its \defin{target marginal} is
\[
 Q_{\mathcal F}(\mathbf x)\coloneqq
 \sum_{u\in\mathcal F(M_1,M_2)}
 \prod_{e\in B(u,0)}x_e.
\]

\begin{proposition}[Lorentzian target marginal]
The polynomial $Q_{\mathcal F}$ is Lorentzian and homogeneous of degree $R$,
or it is zero.
\end{proposition}

\begin{proof}
Apply the basis-generating differential operators of $M_1$ and $M_2$ to the
squarefree monomial $\prod_{e\in E}x_e$. Terms in which the two bases overlap
vanish, and the surviving monomial records the complementary color-$0$ class.
Using the dual basis polynomial, each basis-generating differential operator
can be written as a product of operators $\partial_{x_e}+\partial_{z_e}$
followed by the specialization $\mathbf z=0$. These operations preserve
Lorentzianity by the Brändén--Huh theory \cite{BrandenHuh2020}.
\end{proof}

This Lorentzian statement uses the full choice of the two hidden matroid bases;
it does not require symmetric cycle exchange.

\subsection[rigidWidthThreeTheorem]{The theorem}

\begin{theorem}[Rigid width-three posets]
Let $P$ be a naturally labeled poset with $n$ elements and height $h$. If
$P$ admits a rigid choice, and
\[
 W_P(t)=\sum_{j\geq0}w_jt^j,
\]
then $w_j=0$ for $j>n-h$, and
$(w_0,\dotsc,w_{n-h})$ is ultra-log-concave of order $n-h$ with no internal
zeros.
\end{theorem}

\begin{proof}
The chains $X$ and $Y$ give two transversal matroids on a common set of
positions. A linear extension of $P$ determines a word in the resulting full
three-color family, with target quota $n-h$. The rigidity conditions say that
the interval assigned to a position is independent of the intermediate linear
extension; consequently this map is a bijection, and the number of descents is
the number of marked color-$0$ positions.

After setting the marked variables of $Q_{\mathcal F}$ equal to $t$ and the
remaining variables equal to $s$, we obtain
\[
 \sum_{j=0}^{n-h}w_j s^{n-h-j}t^j.
\]
This polynomial is Lorentzian. The bivariate characterization of Lorentzian
polynomials gives ultra-log-concavity of order $n-h$ and no internal zeros.
\end{proof}

\subsection[leanFormalization]{Lean formalization}

Erik Barkeling supplied a Lean 4 formalization of the theorem and its
supporting construction. The main declarations are:
\begin{itemize}
\item \texttt{ColoredMatroids.WidthThree.isUlc\_of\_natural};
\item \texttt{Choice.isRigid\_of\_closedForm};
\item \texttt{Choice.marginal\_familyU\_mem\_L};
\item \texttt{Choice.isDescentEncoding\_familyU};
\item \texttt{Choice.isUlc\_of\_isRigid}.
\end{itemize}
